\documentclass[11pt]{article}

\usepackage[margin=1in]{geometry}
\usepackage{amsmath, amssymb}
\usepackage{graphicx}
\usepackage{booktabs}
\usepackage{hyperref}
\usepackage{xcolor}
\usepackage{listings}
\usepackage{caption}

\lstset{
  language=Python,
  basicstyle=\ttfamily\small,
  keywordstyle=\color{blue!70!black},
  commentstyle=\color{green!50!black},
  stringstyle=\color{red!60!black},
  breaklines=true,
  frame=single,
  framerule=0.5pt,
  xleftmargin=1em,
  framexleftmargin=0.5em,
  numbers=none,
  showstringspaces=false,
  columns=flexible,
}

\newcommand{\lcq}[1]{\textcolor{red}{#1}}

\title{Learning $\mathbb{Z}_2$ Addition with a Quadratic Autoregressive Model}

\begin{document}
\maketitle

%----------------------------------------------------------------------
\section{Introduction and Data}\label{sec:intro}
%----------------------------------------------------------------------

\subsection{$\mathbb{Z}_2$ addition}

The group $\mathbb{Z}_2 = \{0,1\}$ under addition mod~2 defines the XOR
operation:
\begin{center}
\begin{tabular}{cc|c}
\toprule
$a$ & $b$ & $a \oplus b$ \\
\midrule
0 & 0 & 0 \\
0 & 1 & 1 \\
1 & 0 & 1 \\
1 & 1 & 0 \\
\bottomrule
\end{tabular}
\end{center}

\subsection{Sequence structure}

Sequences have length $T$ (a multiple of~3).  Positions come in
\emph{blocks of~3}: two random bits followed by their XOR.
\begin{equation}\label{eq:block}
  \underbrace{a,\; b,\; a \oplus b}_{\text{block 1}}\;
  \underbrace{a,\; b,\; a \oplus b}_{\text{block 2}}\;
  \cdots
\end{equation}
For example, with $T = 12$:
\[
  0,0,0 \;\mid\; 1,1,0 \;\mid\; 0,1,1 \;\mid\; 1,0,1.
\]

\subsection{True distribution}

Each block is independent.  Within a block the two ``random'' positions
are uniform in $\{0,1\}$ and the third is deterministic, giving exactly
four valid triples per block:
\[
  \{(0,0,0),\;(0,1,1),\;(1,0,1),\;(1,1,0)\},
\]
each with probability $1/4$.  The full sequence probability is
$(1/4)^{T/3}$.

\subsection{Entropy}

Per-position entropy contributions:
\begin{itemize}
  \item Positions $0,1 \pmod{3}$: each contributes $\ln 2$ nats
        (fully random).
  \item Position $2 \pmod{3}$: contributes $0$ nats (deterministic
        given the first two).
\end{itemize}
Average per-position entropy:
\begin{equation}\label{eq:entropy}
  H = \frac{2}{3}\,\ln 2 \approx 0.462\;\text{nats}.
\end{equation}
This is the theoretical minimum loss any correct model can achieve.

%----------------------------------------------------------------------
\section{Why a Linear Model Fails}\label{sec:linear_fails}
%----------------------------------------------------------------------

Suppose we predict $x_2$ from $x_0, x_1$ using a single logistic unit:
\[
  P(x_2 = 1 \mid x_0, x_1) = \sigma(a + b_1 x_0 + b_2 x_1),
\]
where $\sigma(z) = 1/(1 + e^{-z})$.  We need this to match XOR:
\begin{center}
\begin{tabular}{cc|c}
\toprule
$x_0$ & $x_1$ & target $P(x_2=1)$ \\
\midrule
0 & 0 & 0 \\
0 & 1 & 1 \\
1 & 0 & 1 \\
1 & 1 & 0 \\
\bottomrule
\end{tabular}
\end{center}
From rows 1 and~2: increasing $x_1$ (with $x_0=0$) must
\emph{increase} the output, so $b_2 > 0$.  From rows 3 and~4:
increasing $x_1$ (with $x_0=1$) must \emph{decrease} the output, so
$b_2 < 0$.  \textbf{Contradiction.}  No choice of $(a, b_1, b_2)$
works---XOR is not linearly separable.

The accompanying script \texttt{01\_linear\_fails.py} demonstrates this
numerically.

%----------------------------------------------------------------------
\section{Energy-Based vs.\ Autoregressive Quadratic
         Models}\label{sec:subtlety}
%----------------------------------------------------------------------

A natural idea is to use \emph{pairwise (quadratic) interactions} to
capture XOR.  However, \emph{where} the quadratic terms appear matters
fundamentally.

\subsection{Energy-based (Boltzmann) interpretation---doesn't
            work}\label{sec:boltzmann}

Define a joint distribution via an energy function with pairwise terms:
\[
  \log P(\mathbf{x})
  = a + \sum_i b_i x_i + \sum_{i<j} c_{ij}\, x_i x_j - \log Z.
\]
Since $x_t \in \{0,1\}$, the energy splits into terms involving $x_t$
and terms that do not:
\[
  E(\mathbf{x})
  = \underbrace{E_0(\mathbf{x}_{-t})}_{\text{no } x_t}
  + x_t \cdot \underbrace{\bigl(b_t
    + \textstyle\sum_{i \neq t} c_{it}\, x_i\bigr)}_{\displaystyle
    \equiv\;\Delta}.
\]
Computing the conditional:
\[
  P(x_t = 1 \mid \mathbf{x}_{-t})
  = \frac{e^{E_0 + \Delta}/Z}{e^{E_0}/Z + e^{E_0+\Delta}/Z}
  = \frac{e^{\Delta}}{1 + e^{\Delta}}
  = \sigma(\Delta).
\]
Both $Z$ and $e^{E_0}$ cancel.  Crucially,
$\Delta = b_t + \sum_{i\neq t} c_{it}\,x_i$ is \textbf{linear} in the
conditioning variables---the quadratic terms $c_{ij}\,x_i x_j$ where
neither index is~$t$ ended up in $E_0$ and cancelled.  So:
\[
  P(x_t = 1 \mid \text{rest}) = \sigma\!\bigl(b_t
    + \textstyle\sum_{i\neq t} c_{it}\, x_i\bigr),
\]
which is exactly the form proven insufficient for XOR in
Section~\ref{sec:linear_fails}.

\textbf{Key insight:} A pairwise Boltzmann machine has quadratic
interactions in the \emph{joint} but only \emph{linear} conditionals.
Capturing XOR in the conditionals would require a \textbf{three-body}
term $x_0 x_1 x_2$ in the energy---making it a higher-order Boltzmann
machine.

\subsection{Autoregressive interpretation---works}\label{sec:ar}

Instead of an energy-based joint, build the distribution
autoregressively via the chain rule:
\[
  P(x_0, \ldots, x_{T-1})
  = \prod_{t=0}^{T-1} P(x_t \mid x_{<t}),
\]
and put the quadratic features \emph{directly inside each conditional}:
\begin{equation}\label{eq:ar_conditional}
  P(x_t = 1 \mid x_{<t})
  = \sigma\!\Bigl(b_t + \sum_{i<t} W_{ti}\, x_i
    + \sum_{i<j<t} C_{tij}\, x_i x_j\Bigr).
\end{equation}
Now the product $x_i x_j$ appears explicitly in the sigmoid
argument---it does \emph{not} cancel.  This is a fundamentally different
model family from the Boltzmann machine.

\subsection{Concrete solution for position~2}

We need $P(x_2 = 1 \mid x_0, x_1) = x_0 \oplus x_1$.  The logit is:
\[
  \text{logit} = b_2 + W_{20}\,x_0 + W_{21}\,x_1
                 + C_{201}\,x_0\,x_1.
\]
Setting $b_2 = -L$, $W_{20} = W_{21} = 2L$, $C_{201} = -4L$ for
large $L > 0$:
\begin{center}
\begin{tabular}{cc|l|c|c}
\toprule
$x_0$ & $x_1$ & logit & $\sigma(\text{logit})$ & target \\
\midrule
0 & 0 & $-L$         & $\approx 0$ & 0 \\
0 & 1 & $-L+2L = L$  & $\approx 1$ & 1 \\
1 & 0 & $-L+2L = L$  & $\approx 1$ & 1 \\
1 & 1 & $-L+4L-4L = -L$ & $\approx 0$ & 0 \\
\bottomrule
\end{tabular}
\end{center}
As $L \to \infty$ this converges to exact XOR.  The quadratic term
$C_{201}$ is essential---it ``bends'' the decision boundary.

%----------------------------------------------------------------------
\section{The Full Model}\label{sec:full_model}
%----------------------------------------------------------------------

\subsection{The conditional}

Each conditional is parameterised as:
\begin{equation}\label{eq:logit}
  P(x_t = 1 \mid x_{<t}) = \sigma(\Delta_t),
  \qquad
  \Delta_t = b_t + \sum_{i<t} W_{ti}\, x_i
             + \sum_{i<j<t} C_{tij}\, x_i x_j.
\end{equation}
Since $x_t$ is binary, $P(x_t = 0 \mid x_{<t}) = \sigma(-\Delta_t)$.
Both cases are written compactly as
$P(x_t \mid x_{<t}) = \sigma\bigl((2x_t - 1)\,\Delta_t\bigr)$,
since $(2x_t - 1)$ flips the sign of $\Delta_t$ when $x_t = 0$.

\subsection{Chain rule decomposition}

The joint distribution is the product of all conditionals:
\begin{equation}\label{eq:joint}
  P(x_0, \ldots, x_{T-1})
  = \prod_{t=0}^{T-1} \sigma\bigl((2x_t - 1)\,\Delta_t\bigr).
\end{equation}
There is \textbf{no partition function}---each factor is already
normalised ($\sigma(\Delta) + \sigma(-\Delta) = 1$), so the product is
automatically a valid distribution over $\{0,1\}^T$.

\subsection{Computational cost}

Each entry $C_{tij}$ appears in exactly one conditional (the one for
position~$t$, since the mask enforces $i < j < t$).  Computing the full
$\log P(\mathbf{x})$ touches each entry of~$C$ once---total work
$O(T^3)$, the same as a single pass through~$C$.

In practice we need $\log P(\mathbf{x})$ for:
\begin{itemize}
  \item \textbf{Training}: compute $\log P(\mathbf{x})$ for a batch,
        take the gradient, update.  This is what
        \texttt{binary\_cross\_entropy\_with\_logits} computes.
  \item \textbf{Sampling}: go left to right---compute $\Delta_0$,
        sample $x_0 \sim \mathrm{Bernoulli}(\sigma(\Delta_0))$, compute
        $\Delta_1$ (which depends on $x_0$), sample $x_1$, and so on.
        Cost: $O(t^2)$ per position.
\end{itemize}

\subsection{Parameters}

\begin{center}
\begin{tabular}{lllc}
\toprule
Parameter & Shape & Constraint & Count \\
\midrule
$b$       & $(T,)$       & ---              & $T$ \\
$W$       & $(T, T)$     & lower-triangular ($i < t$) & $T(T-1)/2$ \\
$C$       & $(T, T, T)$  & $i < j < t$      & $O(T^3/6)$ \\
\bottomrule
\end{tabular}
\end{center}
For $T = 12$ (4 blocks), $C$ has on the order of $\sim$300 free
parameters.

\subsection{Loss}

Negative log-likelihood per position, averaged over the batch:
\begin{equation}\label{eq:loss}
  \mathcal{L}
  = -\frac{1}{T} \sum_{t=0}^{T-1}
    \mathbb{E}_{\text{data}}\bigl[
      x_t \log\sigma(\Delta_t)
      + (1 - x_t)\log(1 - \sigma(\Delta_t))
    \bigr].
\end{equation}
Theoretical minimum: $(2/3)\ln 2 \approx 0.462$ nats per position.

%----------------------------------------------------------------------
\section{PyTorch Implementation}\label{sec:implementation}
%----------------------------------------------------------------------

The accompanying script \texttt{02\_train\_full.py} implements the full
model.  Below we summarise the key components.

\subsection{Data generation}

Each training batch is generated on the fly (infinite-data regime, no
overfitting).  Block alignment is fixed: every sequence has its XOR
triples at positions $(0,1,2)$, $(3,4,5)$, $(6,7,8)$, etc.  The model
does \emph{not} need to discover where the triples are.  However, the
autoregressive conditionals cross block boundaries---when computing
$P(x_5 \mid x_{<5})$, the model sees positions $0$--$4$, not just the
within-block context $3, 4$.  It must learn that the cross-block weights
(e.g.\ $W_{5,0}$, $C_{5,0,1}$) should be zero.

\subsection{Model class}

The \texttt{QuadraticAutoregressive} module stores three parameter
tensors ($b$, $W$, $C$) and two Boolean masks enforcing causality:
\begin{enumerate}
  \item \textbf{Linear term}: $W$ is masked to be strictly
        lower-triangular, then the linear contribution is computed as
        $\mathbf{x}\, W_{\text{masked}}^\top$.
  \item \textbf{Quadratic term}: the outer product
        $\mathbf{x} \otimes \mathbf{x}$ of shape $(B, T, T)$ is
        contracted against $C_{\text{masked}}$ (with $i < j < t$) via
        \texttt{einsum}, yielding one scalar per (batch, position).
  \item \textbf{Output}: bias $+$ linear $+$ quadratic gives the
        logits (pre-sigmoid values).
\end{enumerate}

\subsection{Training}

Training uses Adam with learning rate $3 \times 10^{-2}$ and batch
size~512.  Fresh data is sampled every step.  Typical convergence for
$T = 12$:
\begin{center}
\begin{tabular}{rl}
\toprule
Step & Loss (nats) \\
\midrule
0    & 0.6931 \\
1000 & 0.4739 \\
2000 & 0.4660 \\
3000 & 0.4635 \\
5000 & 0.4643 \\
\bottomrule
\end{tabular}
\end{center}
The loss converges close to the theoretical minimum of~$0.462$.  Note
that SGD loss is noisy (each step uses a fresh mini-batch), so the
reported loss fluctuates.

\subsection{Inspecting learned parameters}

After training, the parameters exhibit clean structure:
\begin{itemize}
  \item $b_t \approx -L$ at XOR positions, $b_t \approx 0$ at random
        positions.
  \item $W_{t, t-2} \approx W_{t, t-1} \approx 2L$ at XOR positions;
        all other entries near zero.
  \item $C_{t, t-2, t-1} \approx -4L$ at XOR positions; everything else
        near zero.
\end{itemize}
With $L \approx 6.5$, the model independently discovers the XOR solution
for each block.

%----------------------------------------------------------------------
\section{Sparse and Banded Versions}\label{sec:banded}
%----------------------------------------------------------------------

\subsection{Parameter scaling of the full model}

For an order-$n$ interaction term, the source indices must satisfy
$i_1 < i_2 < \cdots < i_n < t$.  The number of such tuples drawn from
$\{0, \ldots, t-1\}$ is $\binom{t}{n}$, so summing over all
positions~$t$:

\begin{center}
\begin{tabular}{llc}
\toprule
Parameter & Total count & Growth \\
\midrule
$b_t$           & $T$                     & $O(T)$     \\
$W_{ti}$        & $T(T-1)/2$              & $O(T^2)$   \\
$C_{tij}$       & $T(T-1)(T-2)/6$         & $O(T^3)$   \\
$D_{tijk}$      & $T(T-1)(T-2)(T-3)/24$   & $O(T^4)$   \\
order-$n$       & $\binom{T}{n+1}$        & $O(T^{n+1})$ \\
\bottomrule
\end{tabular}
\end{center}

For $T = 100$, $C$ alone has $\sim$160k parameters and $D$ has
$\sim$4~million.  This motivates restricting which entries are allowed.

\subsection{Fully banded model---all indices close}

The simplest restriction: position~$t$ can only interact with positions
within a window of size~$k$.  For $C_{tij}$ we require
$t - i \leq k$ and $t - j \leq k$.  Since $i < j < t$, the constraint
$t - i \leq k$ is binding (it is the furthest from~$t$).  This forces
all indices into a window of size~$k$.

\begin{center}
\begin{tabular}{lcc}
\toprule
Parameter & Total & Growth \\
\midrule
$b_t$      & $T$       & $O(T)$      \\
$W_{ti}$   & $kT$      & $O(kT)$     \\
$C_{tij}$  & $k^2 T/2$ & $O(k^2 T)$  \\
$D_{tijk}$ & $k^3 T/6$ & $O(k^3 T)$  \\
order-$n$  & $\binom{k}{n} T$ & $O(k^n T)$ \\
\bottomrule
\end{tabular}
\end{center}

The key feature: \textbf{linear in~$T$} (times a constant depending
on~$k$ and the order).  For the $\mathbb{Z}_2$ data with $k = 3$, $C$
has just $\sim$300 parameters regardless of sequence length.

The accompanying script \texttt{03\_train\_banded.py} implements the
banded model.  With bandwidth~3, it achieves the same loss
($\approx 0.462$) as the full model, confirming purely local structure.

\subsection{Long-range with local pairs}\label{sec:longrange}

The fully banded model assumes all interactions are local: if position~$t$
depends on $x_i x_j$, then both $i$ and~$j$ must be near~$t$.  Some data
may have \emph{long-range} dependencies where distant positions
interact---but the interacting positions are near \emph{each other}.

For $C_{tij}$ the constraint becomes:
\[
  |i - j| \leq k, \qquad i < j < t \quad\text{(causality)},
\]
with no constraint on $t - i$ or $t - j$.

\begin{center}
\begin{tabular}{lcc}
\toprule
Parameter & Total & Growth \\
\midrule
$b_t$      & $T$        & $O(T)$     \\
$W_{ti}$   & $T^2/2$    & $O(T^2)$   \\
$C_{tij}$  & $kT^2/2$   & $O(kT^2)$  \\
$D_{tijk}$ & $k^2 T^2/2$& $O(k^2 T^2)$ \\
order-$n$  & $O(k^{n-1} T^2)$ & $O(k^{n-1} T^2)$ \\
\bottomrule
\end{tabular}
\end{center}

Note that $W$ is unconstrained ($O(T^2)$) because a single source index
has no pair to be ``local'' with.  Banding only kicks in at order~2 and
above.

\subsection{Comparison of scaling}

\begin{center}
\begin{tabular}{lcccc}
\toprule
Scheme & $W$ & $C$ & $D$ & order-$n$ \\
\midrule
Full          & $O(T^2)$ & $O(T^3)$  & $O(T^4)$  & $O(T^{n+1})$ \\
Fully banded  & $O(kT)$  & $O(k^2 T)$& $O(k^3 T)$& $O(k^n T)$ \\
Long-range, local pairs
              & $O(T^2)$ & $O(kT^2)$ & $O(k^2 T^2)$ & $O(k^{n-1} T^2)$ \\
\bottomrule
\end{tabular}
\end{center}

The tradeoff:
\begin{itemize}
  \item \textbf{Fully banded} is linear in~$T$---cheapest, but blind to
        long-range interactions.
  \item \textbf{Long-range, local pairs} is quadratic in~$T$ with one
        fewer power of~$k$.  It lets any position attend to distant
        correlated clusters, as long as the clusters are internally
        compact.
  \item \textbf{Full} is the most expressive but the most expensive,
        with $O(T^{n+1})$ growth.
\end{itemize}

For the $\mathbb{Z}_2$ data, fully banded with $k = 3$ suffices.

\subsubsection{Sparse tensor implementation}

For very large~$T$, storing a dense $(T, T, T)$ tensor is wasteful.
PyTorch sparse COO tensors can materialise only the nonzero entries
of~$C$.  For a banded model with bandwidth~$k$ this is $O(k^2 T)$
entries instead of $O(T^3)$.  At tutorial scale ($T \leq 18$) the
dense-with-mask approach is simpler and fast enough; the sparse approach
pays off at larger~$T$.

%----------------------------------------------------------------------
\section{Exact Training (No Corpus)}\label{sec:exact}
%----------------------------------------------------------------------

The accompanying script \texttt{04\_exact\_training.py} implements this
section.

\subsection{Motivation}

The $\mathbb{Z}_2$ distribution is simple enough that we can
\emph{enumerate} all $2^T$ possible sequences and compute the loss
exactly, eliminating sampling noise.  This is feasible for $T \leq 18$
($2^{18} = 262{,}144$ sequences).

\subsection{True distribution}

Each sequence is either valid (all XOR constraints satisfied, probability
$(1/4)^{T/3}$) or invalid (probability~0).

\subsection{Exact expected loss}

The expected negative log-likelihood under the true distribution is:
\begin{equation}\label{eq:exact_loss}
  \mathcal{L}_{\text{exact}}
  = \sum_{\mathbf{x}} P_{\text{true}}(\mathbf{x})
    \bigl[-\log P_{\text{model}}(\mathbf{x})\bigr],
\end{equation}
where the sum runs over all $2^T$ binary sequences.

\subsection{Training with exact loss}

With $T = 9$ (3~blocks, $2^9 = 512$ sequences), exact training converges
smoothly:
\begin{center}
\begin{tabular}{rl}
\toprule
Step & Exact loss (nats) \\
\midrule
0    & 0.693147 \\
500  & 0.491853 \\
1000 & 0.472786 \\
1500 & 0.467549 \\
2000 & 0.465349 \\
\bottomrule
\end{tabular}
\end{center}
No sampling noise---just clean gradient descent.

\subsection{Comparison: exact vs.\ SGD}

\begin{center}
\begin{tabular}{lcc}
\toprule
Property & SGD on samples & Exact training \\
\midrule
Loss per step & Noisy (mini-batch) & Exact \\
Gradient      & Stochastic estimate & True gradient \\
Cost per step & $O(B \cdot T^2)$   & $O(2^T \cdot T^2)$ \\
Scalability   & Any $T$            & $T \leq {\sim}18$ \\
Final solution& Same               & Same \\
\bottomrule
\end{tabular}
\end{center}

For small~$T$, exact training is a useful debugging tool: if the model
can represent the distribution, exact training will find it without
sampling artefacts.

%----------------------------------------------------------------------
\section{Results and Interpretation}\label{sec:results}
%----------------------------------------------------------------------

\subsection{Learned parameter structure}

After training (either method), the parameters have clean structure:
\begin{itemize}
  \item \textbf{Biases} ($b$): positions $0, 1, 3, 4, 6, 7, \ldots$
        (random positions) have $b \approx 0$; positions
        $2, 5, 8, 11, \ldots$ (XOR positions) have $b \approx -L$ for
        large~$L$.
  \item \textbf{Linear weights} ($W$): XOR position~$t$ depends on
        $t{-}2$ and $t{-}1$ with $W_{t,t-2} \approx W_{t,t-1} \approx 2L$.
        All other entries are near zero.
  \item \textbf{Quadratic weights} ($C$): the only large entries are
        $C_{t,t-2,t-1} \approx -4L$ at XOR positions.
\end{itemize}
The model discovers the \textbf{block-diagonal} structure entirely from
data.

\subsection{Verification checklist}

\begin{enumerate}
  \item \textbf{Loss convergence}: final loss
        $\approx 0.462$ nats/position $= (2/3)\ln 2$. \checkmark
  \item \textbf{XOR prediction accuracy}: sampling 1000 sequences and
        checking predictions at XOR positions gives accuracy
        $\approx 1.0$. \checkmark
  \item \textbf{Banded model}: training with bandwidth~$3$ gives the
        same final loss. \checkmark
  \item \textbf{Exact vs.\ SGD}: both converge to the same parameter
        values (up to optimisation noise). \checkmark
\end{enumerate}

\subsection{Takeaway}

This tutorial demonstrates a key distinction in probabilistic modelling:
\begin{itemize}
  \item \emph{Where} you put interaction terms (in the energy vs.\ in
        the conditionals) determines the expressiveness of your model.
  \item A pairwise energy model has linear conditionals---insufficient
        for XOR.
  \item Putting quadratic features directly in the autoregressive
        conditionals gives the model the capacity to learn XOR.
\end{itemize}

\end{document}
